\documentclass[11pt]{article}
\usepackage[a4paper,margin=1in]{geometry}
\usepackage{amsmath,amssymb,amsthm,mathtools,booktabs}
\usepackage[T1]{fontenc}
\usepackage{lmodern,microtype}
\usepackage[hidelinks]{hyperref}
\newtheorem{theorem}{Theorem}
\newtheorem{lemma}[theorem]{Lemma}
\newtheorem{remark}[theorem]{Remark}
\newcommand{\C}{\mathbb C}
\newcommand{\brk}{\underline{R}}
\DeclareMathOperator{\rank}{rank}
\title{Generic ranks of odd-order Hankel tensors\newline from three compressed slices}
\author{Matthew J. Colbrook\thanks{Department of Applied Mathematics and Theoretical Physics, University of Cambridge, Cambridge, United Kingdom. Email: m.colbrook@damtp.cam.ac.uk.}}
\date{11 September 2026}
\hypersetup{pdfauthor={Matthew J. Colbrook}}
\usepackage{xurl}
\setlength{\emergencystretch}{3em}
\begin{document}
\maketitle
\begin{center}\small Recovered candidate manuscript for TR-13. Not independently refereed.\end{center}
\begin{quote}\small Authorship recorded at the submitter\textquotesingle s request. This recovered draft was generated with AI assistance. The complete original source passed independent agent review against the exact repository target. This is not external human peer review or formal certification; no priority claim is made. Original provenance, full-source review hashes and diagnostic results accompany this manuscript in the submission record.\end{quote}
\begin{abstract}
For every odd order $m\ge3$ and dimension $n\ge2$, the generic ordinary
rank, symmetric rank, ordinary border rank, symmetric border rank and
Vandermonde rank of a complex Hankel tensor all equal
$\lceil(m(n-1)+1)/2\rceil$. The lower bound is certified by an explicit
three-slice Koszul flattening after grouping the tensor factors. Its rank
is established using a two-spike generator and a nilpotent-shift
commutator. This addresses the generic odd-order statement in TR-13; no
claim is made about equality of ranks for every Hankel tensor.
\end{abstract}

\section{Statement and notation}
Use zero-based indices. For $h=(h_0,\ldots,h_D)\in\C^{D+1}$, where
$D=m(n-1)$, write
\[
 H(h)_{i_1,\ldots,i_m}=h_{i_1+\cdots+i_m},
 \qquad 0\le i_j<n.
\]
Let $R$, $R_{\rm sym}$, $\brk$, $\brk_{\rm sym}$ and $R_V$ denote ordinary
rank, symmetric rank, ordinary border rank, symmetric border rank and
Vandermonde rank. A Vandermonde summand is a scalar multiple of
$(1,t,\ldots,t^{n-1})^{\otimes m}$, or its projective counterpart at
infinity. The ordinary border rank is taken in the full tensor space.
The conventions and the earlier even-order and cubic results are in
Nie--Ye \cite{NY}; the remaining generic odd-order assertion is the
canonical TR-13 target \cite{repo}.

\begin{theorem}\label{thm:main}
Fix an odd integer $m\ge3$ and an integer $n\ge2$. There is a nonempty
Zariski-open subset of the Hankel tensor space on which
\[
 R(H)=R_{\rm sym}(H)=\brk(H)=\brk_{\rm sym}(H)=R_V(H)
 =\left\lceil\frac{m(n-1)+1}{2}\right\rceil.
\]
\end{theorem}
Put
\[
 m=2k+1,\quad \ell=n-1,\quad a=k\ell+1,\quad
 s=\lfloor\ell/2\rfloor,\quad r=a+s.
\]
One checks separately for even and odd $\ell$ that
$r=\lceil(D+1)/2\rceil$. The proof consists of a generic Vandermonde upper
bound and a matching ordinary-border-rank lower bound.

\section{An elementary three-slice rank certificate}
\begin{lemma}\label{lem:koszul}
Let $T\in\C^a\otimes\C^3\otimes\C^a$ have slices $M_0,M_1,M_2$ in
its middle factor. Define the $3a\times3a$ matrix
\[
 K(T)=\begin{pmatrix}
 -M_1&M_0&0\\
 -M_2&0&M_0\\
 0&-M_2&M_1
 \end{pmatrix}.
\]
Then $\brk(T)\ge\lceil\rank K(T)/2\rceil$.
If $M_0$ is invertible, then
\[
 \rank K(T)=2a+
 \rank\bigl(M_1M_0^{-1}M_2-M_2M_0^{-1}M_1\bigr).
\]
\end{lemma}
\begin{proof}
For a rank-one tensor, all slices have the form
$M_j=c_j u v^{\mathsf T}$. The displayed block matrix is the tensor
product, up to ordering, of the rank-one matrix $uv^{\mathsf T}$ and
\[
 \begin{pmatrix}-c_1&c_0&0\\-c_2&0&c_0\\0&-c_2&c_1\end{pmatrix},
\]
whose rank is at most two. By subadditivity, $\rank K(T)\le2R(T)$.
Every $(2q+1)$-minor of $K$ is a polynomial in the tensor entries and
vanishes on all tensors of rank at most $q$. It therefore vanishes on
their closure. This proves the border-rank bound without any symmetry
restriction on an approximating sequence.

When $M_0$ is invertible, $(x,y,z)$ belongs to $\ker K$ exactly when
\[
 y=M_0^{-1}M_1x,\qquad z=M_0^{-1}M_2x,
\]
\[
 (M_1M_0^{-1}M_2-M_2M_0^{-1}M_1)x=0.
\]
The formula follows by computing the nullity.
\end{proof}

\section{Compression of the Hankel tensor}
Group the first $k$ factors, the middle factor, and the last $k$ factors.
For every $u\in\{0,\ldots,k\ell\}$ choose once and for all a multi-index
$I(u)\in\{0,\ldots,\ell\}^{k}$ with sum $u$. Such an index exists by
successively allocating at most $\ell$ to each coordinate. Selecting the
coordinates indexed by $I(u)$ is a linear map from the grouped factor
$\C^{n^k}$ onto $\C^a$.

Applying this map to both grouped factors gives a tensor of format
$a\times n\times a$ with entries $h_{u+j+v}$. For $n\ge3$, select the
middle coordinates $0,s,2s$, obtaining $T(h)\in\C^a\otimes\C^3\otimes\C^a$
with slices
\begin{equation}\label{eq:slices}
 (M_j)_{uv}=h_{u+v+js},\qquad 0\le u,v<a,\quad j=0,1,2.
\end{equation}
Grouping and applying fixed linear maps send a rank-one original tensor
to a rank-one tensor or zero. They are also continuous. Hence
\begin{equation}\label{eq:compression}
 \brk(H(h))\ge\brk(T(h)).
\end{equation}
This step does not require a structured decomposition of $H(h)$.

\section{A two-spike generator}
Assume $n\ge3$, so $s\ge1$. Choose
\begin{equation}\label{eq:spikes}
 h_{a-1}=h_{2a+s-1}=1,\qquad h_i=0\text{ for all other }i.
\end{equation}
Both indices lie in $\{0,\ldots,D\}$: the only nontrivial check is
$2a+s-1\le D=2a-2+\ell$, which is equivalent to $s+1\le\ell$.
The indices are distinct. Moreover $a\ge2s+1$, since $k\ge1$ and
$2s\le\ell$.

Let $J$ be the reversal matrix, $J_{uv}=1$ when $u+v=a-1$ and zero
otherwise, and let $L$ be the lower shift $Le_i=e_{i+1}$ for $i<a-1$,
$Le_{a-1}=0$. Define
\[
 E=\sum_{i=0}^{s-1}e_i e_{a-s+i}^{\mathsf T}.
\]
Substitution of \eqref{eq:spikes} into \eqref{eq:slices} gives
\begin{equation}\label{eq:shift}
 M_0=J,\qquad B:=JM_1=L^s,\qquad C:=JM_2=L^{2s}+E.
\end{equation}
For completeness, the first spike in $JM_j$ occurs at column minus row
$-js$, giving $L^{js}$. The second spike cannot occur for $j=0$ or $j=1$;
for $j=2$ it occurs at column minus row $a-s$, giving $E$.

Since powers of $L$ commute, \eqref{eq:shift} yields
\begin{align}\label{eq:comm}
 CB-BC
 &=EL^s-L^sE\nonumber\\
 &=\sum_{i=0}^{s-1}\left(
 e_i e_{a-2s+i}^{\mathsf T}
 -e_{s+i}e_{a-s+i}^{\mathsf T}\right).
\end{align}
In the first $2s$ rows and last $2s$ columns the matrix in
\eqref{eq:comm} is $\operatorname{diag}(I_s,-I_s)$; all other entries
vanish. In particular its rank is exactly $2s$. As $M_0^{-1}=J$,
Lemma~\ref{lem:koszul} gives
\[
 \rank K(T(h))=2a+2s=2r.
\]
Some $2r$-minor of $K(T(h))$ is thus a nonzero polynomial in $h$.
It is nonzero on a nonempty Zariski-open set. On that set,
\eqref{eq:compression} and Lemma~\ref{lem:koszul} imply
\begin{equation}\label{eq:lower}
 \brk(H(h))\ge r.
\end{equation}
The special generator is used only to certify a nonzero minor. It is
not being assumed to lie in the eventual open set on which all five
ranks agree.

If $n=2$, then $s=0$ and $r=a=k+1$. Grouping $k$ factors against $k+1$
factors and selecting one representative for each index sum produces
an $a\times(a+1)$ ordinary matrix flattening. For the generator $h_k=1$
and all other entries zero, its first $a$ columns form $J$. Thus its
generic rank is $a$, proving \eqref{eq:lower} in the binary case as well.

\section{Generic upper bound and conclusion}
Consider the polynomial moment map
\[
 \mathcal M:\C^{2r}\longrightarrow\C^{D+1},\qquad
 ((\lambda_j,t_j)_{j=1}^{r})\longmapsto
 \left(\sum_{j=1}^{r}\lambda_jt_j^i\right)_{i=0}^{D}.
\]
At distinct nodes with every $\lambda_j\ne0$, its Jacobian columns are
$v_D(t_j)$ and $\lambda_jv_D'(t_j)$, where
$v_D(t)=(1,t,\ldots,t^D)^{\mathsf T}$.
A linear functional annihilating all these columns corresponds to a
polynomial $f$ of degree at most $D$ satisfying
$f(t_j)=f'(t_j)=0$ for every $j$.
As $2r\ge D+1$, such a polynomial is zero. The Jacobian therefore has
full row rank $D+1$.

It follows that $\mathcal M$ is dominant. Its image is constructible and
contains a nonempty Zariski-open subset of $\C^{D+1}$; equivalently one
may apply the standard dominance and constructibility facts for
polynomial maps. For every $h$ in this image,
\[
 H(h)=\sum_{j=1}^{r}\lambda_j
       (1,t_j,\ldots,t_j^{\ell})^{\otimes m},
\]
so $R_V(H(h))\le r$.
Intersect this nonempty open set with the nonempty open set from
\eqref{eq:lower}. Their intersection is nonempty because the Hankel
parameter space is irreducible. On the intersection, the inequalities
\[
 \brk\le R\le R_{\rm sym}\le R_V,
 \qquad
 \brk\le\brk_{\rm sym}\le R_{\rm sym}
\]
force every displayed rank to equal $r$. This proves
Theorem~\ref{thm:main}.

\section{Scope and reproducibility}
Theorem~\ref{thm:main} addresses the full generic statement for every
admissible pair $(m,n)$ in TR-13. It also supplies a second proof for
order three, which was already known \cite{NY}. It is not a proof of
Comon's conjecture for all Hankel tensors (TR-14): the nonzero-minor and
moment-map arguments both specify open sets and do not eliminate
exceptional tensors.

The accompanying exact-arithmetic script \texttt{code/verify\_tr13.py}
checks the slice identities, the entire commutator, and its indicated
nonzero minor over a parameter grid. For small formats it independently
builds the full Koszul matrix and computes its rank, and verifies that
rank-one inputs have Koszul rank at most two. No floating-point rank
tolerance enters these checks. The parameter-uniform proof is the one
above; the grid is only a diagnostic.

\begin{thebibliography}{9}
\bibitem{NY}
J. Nie and K. Ye, \emph{Hankel Tensor Decompositions and Ranks},
SIAM Journal on Matrix Analysis and Applications 40 (2019),
\href{https://doi.org/10.1137/18M1168285}{doi:10.1137/18M1168285}.
Author preprint, Sections 4--6, especially Question 6.1 and Conjecture 6.2:
\url{https://arxiv.org/abs/1706.03631}.
\bibitem{repo}
OpenProblemsInNLA, \emph{TR-13: Equality of ranks for generic odd-order
Hankel tensors}, canonical statement.
\url{https://github.com/ajt60gaibb/OpenProblemsInNLA/tree/main/tensor-computations/TR-13}.
\end{thebibliography}
\end{document}
